\IfFileExists{stacks-project.cls}{%
\documentclass{stacks-project}
}{%
\documentclass{amsart}
}

% For dealing with references we use the comment environment
\usepackage{verbatim}
\newenvironment{reference}{\comment}{\endcomment}
%\newenvironment{reference}{}{}
\newenvironment{slogan}{\comment}{\endcomment}
\newenvironment{history}{\comment}{\endcomment}

% For commutative diagrams we use Xy-pic
\usepackage[all]{xy}

% We use 2cell for 2-commutative diagrams.
\xyoption{2cell}
\UseAllTwocells

% We use multicol for the list of chapters between chapters
\usepackage{multicol}

% This is generally recommended for better output
\usepackage{lmodern}
\usepackage[T1]{fontenc}

% For cross-file-references
\usepackage{xr-hyper}

% Package for hypertext links:
\usepackage{hyperref}

% For any local file, say "hello.tex" you want to link to please
% use \externaldocument[hello-]{hello}
\externaldocument[introduction-]{introduction}
\externaldocument[conventions-]{conventions}
\externaldocument[sets-]{sets}
\externaldocument[categories-]{categories}
\externaldocument[topology-]{topology}
\externaldocument[sheaves-]{sheaves}
\externaldocument[sites-]{sites}
\externaldocument[stacks-]{stacks}
\externaldocument[fields-]{fields}
\externaldocument[algebra-]{algebra}
\externaldocument[brauer-]{brauer}
\externaldocument[homology-]{homology}
\externaldocument[derived-]{derived}
\externaldocument[simplicial-]{simplicial}
\externaldocument[more-algebra-]{more-algebra}
\externaldocument[smoothing-]{smoothing}
\externaldocument[modules-]{modules}
\externaldocument[sites-modules-]{sites-modules}
\externaldocument[injectives-]{injectives}
\externaldocument[cohomology-]{cohomology}
\externaldocument[sites-cohomology-]{sites-cohomology}
\externaldocument[dga-]{dga}
\externaldocument[dpa-]{dpa}
\externaldocument[sdga-]{sdga}
\externaldocument[hypercovering-]{hypercovering}
\externaldocument[schemes-]{schemes}
\externaldocument[constructions-]{constructions}
\externaldocument[properties-]{properties}
\externaldocument[morphisms-]{morphisms}
\externaldocument[coherent-]{coherent}
\externaldocument[divisors-]{divisors}
\externaldocument[limits-]{limits}
\externaldocument[varieties-]{varieties}
\externaldocument[topologies-]{topologies}
\externaldocument[descent-]{descent}
\externaldocument[perfect-]{perfect}
\externaldocument[more-morphisms-]{more-morphisms}
\externaldocument[flat-]{flat}
\externaldocument[groupoids-]{groupoids}
\externaldocument[more-groupoids-]{more-groupoids}
\externaldocument[etale-]{etale}
\externaldocument[chow-]{chow}
\externaldocument[intersection-]{intersection}
\externaldocument[pic-]{pic}
\externaldocument[weil-]{weil}
\externaldocument[adequate-]{adequate}
\externaldocument[dualizing-]{dualizing}
\externaldocument[duality-]{duality}
\externaldocument[discriminant-]{discriminant}
\externaldocument[derham-]{derham}
\externaldocument[local-cohomology-]{local-cohomology}
\externaldocument[algebraization-]{algebraization}
\externaldocument[curves-]{curves}
\externaldocument[resolve-]{resolve}
\externaldocument[models-]{models}
\externaldocument[functors-]{functors}
\externaldocument[equiv-]{equiv}
\externaldocument[pione-]{pione}
\externaldocument[etale-cohomology-]{etale-cohomology}
\externaldocument[proetale-]{proetale}
\externaldocument[relative-cycles-]{relative-cycles}
\externaldocument[more-etale-]{more-etale}
\externaldocument[trace-]{trace}
\externaldocument[crystalline-]{crystalline}
\externaldocument[spaces-]{spaces}
\externaldocument[spaces-properties-]{spaces-properties}
\externaldocument[spaces-morphisms-]{spaces-morphisms}
\externaldocument[decent-spaces-]{decent-spaces}
\externaldocument[spaces-cohomology-]{spaces-cohomology}
\externaldocument[spaces-limits-]{spaces-limits}
\externaldocument[spaces-divisors-]{spaces-divisors}
\externaldocument[spaces-over-fields-]{spaces-over-fields}
\externaldocument[spaces-topologies-]{spaces-topologies}
\externaldocument[spaces-descent-]{spaces-descent}
\externaldocument[spaces-perfect-]{spaces-perfect}
\externaldocument[spaces-more-morphisms-]{spaces-more-morphisms}
\externaldocument[spaces-flat-]{spaces-flat}
\externaldocument[spaces-groupoids-]{spaces-groupoids}
\externaldocument[spaces-more-groupoids-]{spaces-more-groupoids}
\externaldocument[bootstrap-]{bootstrap}
\externaldocument[spaces-pushouts-]{spaces-pushouts}
\externaldocument[spaces-chow-]{spaces-chow}
\externaldocument[groupoids-quotients-]{groupoids-quotients}
\externaldocument[spaces-more-cohomology-]{spaces-more-cohomology}
\externaldocument[spaces-simplicial-]{spaces-simplicial}
\externaldocument[spaces-duality-]{spaces-duality}
\externaldocument[formal-spaces-]{formal-spaces}
\externaldocument[restricted-]{restricted}
\externaldocument[spaces-resolve-]{spaces-resolve}
\externaldocument[formal-defos-]{formal-defos}
\externaldocument[defos-]{defos}
\externaldocument[cotangent-]{cotangent}
\externaldocument[examples-defos-]{examples-defos}
\externaldocument[algebraic-]{algebraic}
\externaldocument[examples-stacks-]{examples-stacks}
\externaldocument[stacks-sheaves-]{stacks-sheaves}
\externaldocument[criteria-]{criteria}
\externaldocument[artin-]{artin}
\externaldocument[quot-]{quot}
\externaldocument[stacks-properties-]{stacks-properties}
\externaldocument[stacks-morphisms-]{stacks-morphisms}
\externaldocument[stacks-limits-]{stacks-limits}
\externaldocument[stacks-cohomology-]{stacks-cohomology}
\externaldocument[stacks-perfect-]{stacks-perfect}
\externaldocument[stacks-introduction-]{stacks-introduction}
\externaldocument[stacks-more-morphisms-]{stacks-more-morphisms}
\externaldocument[stacks-geometry-]{stacks-geometry}
\externaldocument[moduli-]{moduli}
\externaldocument[moduli-curves-]{moduli-curves}
\externaldocument[examples-]{examples}
\externaldocument[exercises-]{exercises}
\externaldocument[guide-]{guide}
\externaldocument[desirables-]{desirables}
\externaldocument[coding-]{coding}
\externaldocument[obsolete-]{obsolete}
\externaldocument[fdl-]{fdl}
\externaldocument[index-]{index}

% Theorem environments.
%
\theoremstyle{plain}
\newtheorem{theorem}[subsection]{Theorem}
\newtheorem{proposition}[subsection]{Proposition}
\newtheorem{lemma}[subsection]{Lemma}

\theoremstyle{definition}
\newtheorem{definition}[subsection]{Definition}
\newtheorem{example}[subsection]{Example}
\newtheorem{exercise}[subsection]{Exercise}
\newtheorem{situation}[subsection]{Situation}

\theoremstyle{remark}
\newtheorem{remark}[subsection]{Remark}
\newtheorem{remarks}[subsection]{Remarks}

\numberwithin{equation}{subsection}

% Macros
%
\def\lim{\mathop{\mathrm{lim}}\nolimits}
\def\colim{\mathop{\mathrm{colim}}\nolimits}
\def\Spec{\mathop{\mathrm{Spec}}}
\def\Hom{\mathop{\mathrm{Hom}}\nolimits}
\def\Ext{\mathop{\mathrm{Ext}}\nolimits}
\def\SheafHom{\mathop{\mathcal{H}\!\mathit{om}}\nolimits}
\def\SheafExt{\mathop{\mathcal{E}\!\mathit{xt}}\nolimits}
\def\Sch{\mathit{Sch}}
\def\Mor{\mathop{\mathrm{Mor}}\nolimits}
\def\Ob{\mathop{\mathrm{Ob}}\nolimits}
\def\Sh{\mathop{\mathit{Sh}}\nolimits}
\def\NL{\mathop{N\!L}\nolimits}
\def\CH{\mathop{\mathrm{CH}}\nolimits}
\def\proetale{{pro\text{-}\acute{e}tale}}
\def\etale{{\acute{e}tale}}
\def\QCoh{\mathit{QCoh}}
\def\Ker{\mathop{\mathrm{Ker}}}
\def\Im{\mathop{\mathrm{Im}}}
\def\Coker{\mathop{\mathrm{Coker}}}
\def\Coim{\mathop{\mathrm{Coim}}}

% Boxtimes
%
\DeclareMathSymbol{\boxtimes}{\mathbin}{AMSa}{"02}

%
% Macros for moduli stacks/spaces
%
\def\QCohstack{\mathcal{QC}\!\mathit{oh}}
\def\Cohstack{\mathcal{C}\!\mathit{oh}}
\def\Spacesstack{\mathcal{S}\!\mathit{paces}}
\def\Quotfunctor{\mathrm{Quot}}
\def\Hilbfunctor{\mathrm{Hilb}}
\def\Curvesstack{\mathcal{C}\!\mathit{urves}}
\def\Polarizedstack{\mathcal{P}\!\mathit{olarized}}
\def\Complexesstack{\mathcal{C}\!\mathit{omplexes}}
% \Pic is the operator that assigns to X its picard group, usage \Pic(X)
% \Picardstack_{X/B} denotes the Picard stack of X over B
% \Picardfunctor_{X/B} denotes the Picard functor of X over B
\def\Pic{\mathop{\mathrm{Pic}}\nolimits}
\def\Picardstack{\mathcal{P}\!\mathit{ic}}
\def\Picardfunctor{\mathrm{Pic}}
\def\Deformationcategory{\mathcal{D}\!\mathit{ef}}
\usepackage{cleveref}
\crefname{originalreference}{original reference}{original references}
\Crefname{originalreference}{Original reference}{Original references}
\makeatletter
\newlabel{decent-spaces-lemma-decent-space-elementary-etale-neighbourhood}{{0BBP}{source}{}{sourcepointer-1}{}}
\newlabel{decent-spaces-lemma-decent-space-elementary-etale-neighbourhood@cref}{{[originalreference][0][]0BBP}{[1][0][]source}{}{}{}}
\makeatother
\crefname{originallemma}{original lemma}{original lemmas}
\Crefname{originallemma}{Original lemma}{Original lemmas}
\makeatletter
\newlabel{lemma-equivalence}{{0BHB}{source}{}{sourcepointer-2}{}}
\newlabel{lemma-equivalence@cref}{{[originallemma][0][]0BHB}{[1][0][]source}{}{}{}}
\makeatother
\crefname{originallemma}{original lemma}{original lemmas}
\Crefname{originallemma}{Original lemma}{Original lemmas}
\makeatletter
\newlabel{lemma-equivalence-sequence-normalized-blowups}{{0BHL}{source}{}{sourcepointer-3}{}}
\newlabel{lemma-equivalence-sequence-normalized-blowups@cref}{{[originallemma][0][]0BHL}{[1][0][]source}{}{}{}}
\makeatother
\crefname{originalreference}{original reference}{original references}
\Crefname{originalreference}{Original reference}{Original references}
\makeatletter
\newlabel{resolve-lemma-dominate-by-normalized-blowing-up}{{0BHM}{source}{}{sourcepointer-4}{}}
\newlabel{resolve-lemma-dominate-by-normalized-blowing-up@cref}{{[originalreference][0][]0BHM}{[1][0][]source}{}{}{}}
\makeatother
\crefname{originalsection}{original section}{original sections}
\Crefname{originalsection}{Original section}{Original sections}
\makeatletter
\newlabel{section-strategy}{{0BH9}{source}{}{sourcepointer-5}{}}
\newlabel{section-strategy@cref}{{[originalsection][0][]0BH9}{[1][0][]source}{}{}{}}
\makeatother
\crefname{originalreference}{original reference}{original references}
\Crefname{originalreference}{Original reference}{Original references}
\makeatletter
\newlabel{spaces-more-morphisms-lemma-dominate-modification-by-blowup}{{087G}{source}{}{sourcepointer-6}{}}
\newlabel{spaces-more-morphisms-lemma-dominate-modification-by-blowup@cref}{{[originalreference][0][]087G}{[1][0][]source}{}{}{}}
\makeatother
\crefname{originalreference}{original reference}{original references}
\Crefname{originalreference}{Original reference}{Original references}
\makeatletter
\newlabel{spaces-morphisms-section-dimension-formula}{{0BAW}{source}{}{sourcepointer-7}{}}
\newlabel{spaces-morphisms-section-dimension-formula@cref}{{[originalreference][0][]0BAW}{[1][0][]source}{}{}{}}
\makeatother
\crefname{originalreference}{original reference}{original references}
\Crefname{originalreference}{Original reference}{Original references}
\makeatletter
\newlabel{spaces-morphisms-section-nagata}{{0BAT}{source}{}{sourcepointer-8}{}}
\newlabel{spaces-morphisms-section-nagata@cref}{{[originalreference][0][]0BAT}{[1][0][]source}{}{}{}}
\makeatother
\crefname{originalreference}{original reference}{original references}
\Crefname{originalreference}{Original reference}{Original references}
\makeatletter
\newlabel{spaces-morphisms-section-normalization}{{07U3}{source}{}{sourcepointer-9}{}}
\newlabel{spaces-morphisms-section-normalization@cref}{{[originalreference][0][]07U3}{[1][0][]source}{}{}{}}
\makeatother
\crefname{originalreference}{original reference}{original references}
\Crefname{originalreference}{Original reference}{Original references}
\makeatletter
\newlabel{spaces-over-fields-lemma-modification-normal-iso-over-codimension-1}{{0BBQ}{source}{}{sourcepointer-10}{}}
\newlabel{spaces-over-fields-lemma-modification-normal-iso-over-codimension-1@cref}{{[originalreference][0][]0BBQ}{[1][0][]source}{}{}{}}
\makeatother
\title{A proper birational morphism of 2-dimensional Nagata spaces is dominated by a sequence of normalized blowups in closed points}
\author{The Stacks Project Authors}
\date{stacks-project:spaces-resolve.tex:local-snapshot-20261012; GNU FDL 1.2 or later}
\begin{document}
\maketitle
\noindent\textit{Editorial note.} Two same-chapter inputs are NOT carded and are declared external: lemma-equivalence-sequence-normalized-blowups (zero-content cross-reference to the blowup-sequence proof) and lemma-equivalence (dropped: fully-faithfulness part explicitly omitted).\par\medskip
\noindent\textit{Allowed standard prerequisites (editorial).}\begin{itemize}
\item Normalizations are normal and morphisms factor through them. Morphisms of Spaces, section normalization (cross-chapter, named). Lets the proof replace X and Y by their normalizations.
\item Modifications of normal spaces are isomorphisms off a finite set of closed points. Spaces over Fields, lemma modification-normal-iso-over-codimension-1 (cross-chapter, named). Produces the finite set T with f an isomorphism over X minus T.
\item Dominating modifications by admissible blowups. More on Morphisms of Spaces, lemma dominate-modification-by-blowup (cross-chapter, named). Gives an X minus T admissible blowup Y' -> X dominating Y, made representable after normalization.
\item Scheme-case domination by normalized blowups. Resolution of Surfaces, lemma dominate-by-normalized-blowing-up (cross-chapter, named). Provides the normalized blowup sequence on each etale neighbourhood dominating Y\_i.
\item Normalized blowup sequences under the etale neighbourhood equivalence. lemma-equivalence-sequence-normalized-blowups (same chapter, NOT carded, declared external). Assembles the per-neighbourhood sequences into one global sequence.
\item The etale neighbourhood functor is an equivalence. lemma-equivalence (same chapter, NOT carded, declared external). Concludes X\_m dominates Y from the per-neighbourhood domination.
\end{itemize}\medskip
\ifdefined\hypertarget\hypertarget{sourcepointer-1}{}\fi
\noindent\textit{Source reference.} Original reference 0BBP: focused-seeds/geometry/S9-spres-dominate-by-normalized-blowing-up/ref-tags.txt, source lines 5--5. Navigation pointer only: the counted proof cites this statement by \textbackslash{}ref as a standing input; the deduction the counted proof has to make is carried out inside the counted span, and no part of that deduction is delegated to this pointer.\par\medskip
\ifdefined\hypertarget\hypertarget{sourcepointer-2}{}\fi
\noindent\textit{Source reference.} Original lemma 0BHB: sources/stacks-spaces-resolve/pinned-original-source.tex, source lines 425--425. Navigation pointer only: the counted proof cites this statement by \textbackslash{}ref as a standing input; the deduction the counted proof has to make is carried out inside the counted span, and no part of that deduction is delegated to this pointer.\par\medskip
\ifdefined\hypertarget\hypertarget{sourcepointer-3}{}\fi
\noindent\textit{Source reference.} Original lemma 0BHL: sources/stacks-spaces-resolve/pinned-original-source.tex, source lines 708--708. Navigation pointer only: the counted proof cites this statement by \textbackslash{}ref as a standing input; the deduction the counted proof has to make is carried out inside the counted span, and no part of that deduction is delegated to this pointer.\par\medskip
\ifdefined\hypertarget\hypertarget{sourcepointer-4}{}\fi
\noindent\textit{Source reference.} Original reference 0BHM: focused-seeds/geometry/S9-spres-dominate-by-normalized-blowing-up/ref-tags.txt, source lines 6--6. Navigation pointer only: the counted proof cites this statement by \textbackslash{}ref as a standing input; the deduction the counted proof has to make is carried out inside the counted span, and no part of that deduction is delegated to this pointer.\par\medskip
\ifdefined\hypertarget\hypertarget{sourcepointer-5}{}\fi
\noindent\textit{Source reference.} Original section 0BH9: sources/stacks-spaces-resolve/pinned-original-source.tex, source lines 391--391. Navigation pointer only: the counted proof cites this statement by \textbackslash{}ref as a standing input; the deduction the counted proof has to make is carried out inside the counted span, and no part of that deduction is delegated to this pointer.\par\medskip
\ifdefined\hypertarget\hypertarget{sourcepointer-6}{}\fi
\noindent\textit{Source reference.} Original reference 087G: focused-seeds/geometry/S9-spres-dominate-by-normalized-blowing-up/ref-tags.txt, source lines 7--7. Navigation pointer only: the counted proof cites this statement by \textbackslash{}ref as a standing input; the deduction the counted proof has to make is carried out inside the counted span, and no part of that deduction is delegated to this pointer.\par\medskip
\ifdefined\hypertarget\hypertarget{sourcepointer-7}{}\fi
\noindent\textit{Source reference.} Original reference 0BAW: focused-seeds/geometry/S9-spres-dominate-by-normalized-blowing-up/ref-tags.txt, source lines 8--8. Navigation pointer only: the counted proof cites this statement by \textbackslash{}ref as a standing input; the deduction the counted proof has to make is carried out inside the counted span, and no part of that deduction is delegated to this pointer.\par\medskip
\ifdefined\hypertarget\hypertarget{sourcepointer-8}{}\fi
\noindent\textit{Source reference.} Original reference 0BAT: focused-seeds/geometry/S9-spres-dominate-by-normalized-blowing-up/ref-tags.txt, source lines 9--9. Navigation pointer only: the counted proof cites this statement by \textbackslash{}ref as a standing input; the deduction the counted proof has to make is carried out inside the counted span, and no part of that deduction is delegated to this pointer.\par\medskip
\ifdefined\hypertarget\hypertarget{sourcepointer-9}{}\fi
\noindent\textit{Source reference.} Original reference 07U3: focused-seeds/geometry/S9-spres-dominate-by-normalized-blowing-up/ref-tags.txt, source lines 10--10. Navigation pointer only: the counted proof cites this statement by \textbackslash{}ref as a standing input; the deduction the counted proof has to make is carried out inside the counted span, and no part of that deduction is delegated to this pointer.\par\medskip
\ifdefined\hypertarget\hypertarget{sourcepointer-10}{}\fi
\noindent\textit{Source reference.} Original reference 0BBQ: focused-seeds/geometry/S9-spres-dominate-by-normalized-blowing-up/ref-tags.txt, source lines 11--11. Navigation pointer only: the counted proof cites this statement by \textbackslash{}ref as a standing input; the deduction the counted proof has to make is carried out inside the counted span, and no part of that deduction is delegated to this pointer.\par\medskip
\begin{lemma}
\label{lemma-dominate-by-normalized-blowing-up}
Let $S$ be a scheme. Let $X$ be a Noetherian Nagata algebraic space over $S$
with $\dim(X) = 2$. Let $f : Y \to X$ be a proper birational morphism.
Then there exists a commutative diagram
$$
\xymatrix{
X_n \ar[r] \ar[d] &
X_{n - 1} \ar[r] &
\ldots \ar[r] &
X_1 \ar[r] &
X_0 \ar[d] \\
Y \ar[rrrr]  & & & & X
}
$$
where $X_0 \to X$ is the normalization and
where $X_{i + 1} \to X_i$ is the normalized blowing up of $X_i$ at a closed
point.
\end{lemma}

\begin{proof}
Although one can prove this lemma directly for algebraic spaces,
we will continue the approach used above to reduce it to the case
of schemes.

\medskip\noindent
We will use that Noetherian algebraic spaces are quasi-separated
and hence points have well defined residue fields (for example by
Decent Spaces, Lemma
\ref{decent-spaces-lemma-decent-space-elementary-etale-neighbourhood}).
We will use the results of Morphisms of Spaces, Sections
\ref{spaces-morphisms-section-nagata},
\ref{spaces-morphisms-section-dimension-formula}, and
\ref{spaces-morphisms-section-normalization} without further mention.
We may replace $Y$ by its normalization. Let $X_0 \to X$ be the normalization.
The morphism $Y \to X$ factors through $X_0$.
Thus we may assume that both $X$ and $Y$ are normal.

\medskip\noindent
Assume $X$ and $Y$ are normal. The morphism $f : Y \to X$ is an isomorphism
over an open which contains every point of codimension $0$ and $1$ in $Y$ and
every point of $Y$ over which the fibre is finite, see
Spaces over Fields, Lemma
\ref{spaces-over-fields-lemma-modification-normal-iso-over-codimension-1}.
Hence we see that there is a finite set of closed points $T \subset |X|$
such that $f$ is an isomorphism over $X \setminus T$.
By More on Morphisms of Spaces, Lemma
\ref{spaces-more-morphisms-lemma-dominate-modification-by-blowup}
there exists an $X \setminus T$-admissible blowup $Y' \to X$
which dominates $Y$. After replacing $Y$ by the normalization of
$Y'$ we see that we may assume that $Y \to X$ is representable.

\medskip\noindent
Say $T = \{x_1, \ldots, x_r\}$. Pick elementary \'etale neighbourhoods
$(U_i, u_i) \to (X, x_i)$ as in Section \ref{section-strategy}.
For each $i$ the morphism $Y_i = Y \times_X U_i \to U_i$
is a proper birational morphism which is an isomorphism over
$U_i \setminus \{u_i\}$. Thus we may apply
Resolution of Surfaces, Lemma
\ref{resolve-lemma-dominate-by-normalized-blowing-up}
to find a sequence
$$
X_{i, m_i} \to X_{i, m_i - 1} \to \ldots \to X_1 \to X_{i, 0} = U_i
$$
of normalized blowing ups in closed points lying over $u_i$ such that
$X_{i, m_i}$ dominates $Y_i$.
By Lemma \ref{lemma-equivalence-sequence-normalized-blowups}
we find a sequence of normalized blowing ups
$$
X_m \to X_{m - 1} \to \ldots \to X_1 \to X_0 = X
$$
as in the statement of the lemma whose base change to our $U_i$
produces the given sequences. It follows that $X_m$ dominates
$Y$ by the equivalence of categories of
Lemma \ref{lemma-equivalence}.
\end{proof}

\end{document}
